\documentclass[12pt,a4paper]{ctexart}
\usepackage[left=2.5cm,right=2.5cm,top=2.3cm,bottom=2.3cm]{geometry}
\usepackage{xcolor}
\usepackage{tcolorbox}
\usepackage{titlesec}
\usepackage{fancyhdr}
\usepackage{enumitem}
\usepackage{tabularx}
\usepackage{amsmath,amssymb}
\usepackage{hyperref}
\tcbuselibrary{skins,breakable}

\definecolor{MainColor}{HTML}{2A6F73}
\definecolor{SecondColor}{HTML}{3CAEA3}
\definecolor{LightMain}{HTML}{EAF7F6}
\definecolor{TitleDark}{HTML}{1F3A40}
\definecolor{TextGray}{HTML}{555555}
\definecolor{BorderGray}{HTML}{CBD5D6}
\definecolor{WarnColor}{HTML}{F2B35D}
\definecolor{ErrorColor}{HTML}{C94C4C}
\definecolor{DefineColor}{HTML}{6C63B7}

\linespread{1.25}
\setlength{\parindent}{2em}
\setlength{\parskip}{0.25em}
\hypersetup{colorlinks=true,linkcolor=MainColor,urlcolor=MainColor}
\pagestyle{fancy}
\fancyhf{}
\fancyfoot[L]{\textcolor{MainColor}{机器学习笔记}}
\fancyfoot[R]{\textcolor{MainColor}{\thepage}}
\renewcommand{\headrulewidth}{0pt}
\renewcommand{\footrulewidth}{0pt}

\newcommand{\topbar}[2]{\noindent{\small\textcolor{MainColor}{#1}}\hfill{\small\bfseries\textcolor{TitleDark}{#2}}\par\vspace{0.35em}{\color{BorderGray}\hrule}\vspace{1.2em}}
\newcommand{\notetitlecard}[6]{
\begin{tcolorbox}[enhanced,colback=white,colframe=BorderGray,boxrule=0.6pt,arc=2mm,left=12pt,right=12pt,top=10pt,bottom=10pt,borderline west={4pt}{0pt}{SecondColor}]
\begin{tabularx}{\linewidth}{X r}
{\small\textcolor{TextGray}{#1}} & {\small\textcolor{TextGray}{#2}} \\[0.9em]
\multicolumn{2}{l}{{\Huge\bfseries\textcolor{MainColor}{#3}}} \\[0.45em]
\multicolumn{2}{l}{{\large\textcolor{SecondColor!90!black}{#4}}} \\[0.45em]
\multicolumn{2}{l}{{\small\textcolor{TextGray}{课程：#5 \qquad 作者：#6}}}
\end{tabularx}\end{tcolorbox}\vspace{1em}}
\newtcolorbox{summarybox}[1][本节摘要]{enhanced,breakable,colback=LightMain,colframe=SecondColor,colbacktitle=SecondColor,coltitle=white,title={#1},fonttitle=\bfseries,boxrule=0.6pt,arc=2mm,left=12pt,right=12pt,top=10pt,bottom=10pt}
\newtcolorbox{definitionbox}[1][定义]{enhanced,breakable,colback=DefineColor!6,colframe=DefineColor,colbacktitle=DefineColor,coltitle=white,title={#1},fonttitle=\bfseries,boxrule=0.6pt,arc=2mm,left=12pt,right=12pt,top=10pt,bottom=10pt}
\newtcolorbox{keybox}[1][重点]{enhanced,breakable,colback=MainColor!5,colframe=MainColor,colbacktitle=MainColor,coltitle=white,title={#1},fonttitle=\bfseries,boxrule=0.6pt,arc=2mm,left=12pt,right=12pt,top=10pt,bottom=10pt}
\newtcolorbox{examplebox}[1][例题]{enhanced,breakable,colback=white,colframe=SecondColor,colbacktitle=SecondColor!85!black,coltitle=white,title={#1},fonttitle=\bfseries,boxrule=0.6pt,arc=2mm,left=12pt,right=12pt,top=10pt,bottom=10pt}
\newtcolorbox{mistakebox}[1][易错点]{enhanced,breakable,colback=ErrorColor!5,colframe=ErrorColor,colbacktitle=ErrorColor,coltitle=white,title={#1},fonttitle=\bfseries,boxrule=0.6pt,arc=2mm,left=12pt,right=12pt,top=10pt,bottom=10pt}
\newtcolorbox{tipbox}[1][提示]{enhanced,breakable,colback=WarnColor!10,colframe=WarnColor,colbacktitle=WarnColor,coltitle=white,title={#1},fonttitle=\bfseries,boxrule=0.6pt,arc=2mm,left=12pt,right=12pt,top=10pt,bottom=10pt}
\titleformat{\section}{\Large\bfseries\color{MainColor}}{\thesection}{0.8em}{}
\titleformat{\subsection}{\large\bfseries\color{TitleDark}}{\thesubsection}{0.8em}{}
\titleformat{\subsubsection}{\normalsize\bfseries\color{SecondColor!80!black}}{\thesubsubsection}{0.8em}{}
\setlist[itemize]{leftmargin=2.2em,itemsep=0.3em,topsep=0.4em}
\setlist[enumerate]{leftmargin=2.2em,itemsep=0.3em,topsep=0.4em}
\newcommand{\important}[1]{\textbf{\textcolor{MainColor}{#1}}}
\newcommand{\warning}[1]{\textbf{\textcolor{ErrorColor}{#1}}}
\newcommand{\concept}[1]{\textbf{\textcolor{DefineColor}{#1}}}

\begin{document}
\topbar{吴恩达机器学习}{学习笔记}
\notetitlecard{Study Notes Series}{2026 Spring}{吴恩达机器学习笔记(8)}{}{机器学习}{C++与草稿纸}

\begin{summarybox}
本周回到无监督学习，介绍 K-均值聚类和主成分分析（PCA）。K-均值把样本分配给若干中心，PCA 则寻找低维方向以压缩和可视化数据。两者都需要明确目标、尺度和结果解释。
\end{summarybox}

\section{K-均值聚类}
\begin{definitionbox}[无监督学习]
无监督学习没有给定标签 $y$，算法只根据输入样本 $x^{(i)}$ 寻找数据中的结构。聚类的目标是把相似样本放在同一组，把不同样本分开。
\end{definitionbox}
\begin{keybox}[K-均值算法]
给定簇数 $K$，K-均值重复两个步骤：
\begin{enumerate}
  \item \textbf{分配：}将每个样本分给距离最近的聚类中心 $\mu_k$；
  \item \textbf{移动：}把每个中心更新为所属样本的均值。
\end{enumerate}
直到样本分配不再变化或目标函数改善很小。算法输出的是簇编号和中心，不是预先存在的类别名称。
\end{keybox}
\begin{definitionbox}[优化目标]
K-均值最小化簇内平方距离：
\[
J(c,\mu)=\frac1m\sum_{i=1}^{m}\|x^{(i)}-\mu_{c^{(i)}}\|^2,
\]
其中 $c^{(i)}$ 是样本所属簇。分配步骤对固定中心最小化 $c$，移动步骤对固定分配最小化 $\mu$，所以每次迭代不会增加目标函数。
\end{definitionbox}
\begin{tipbox}[随机初始化与局部最优]
K-均值的目标通常非凸，可能收敛到局部最优。随机选择不同初始中心并重复运行，保留目标函数最小的结果。若 $K$ 较小，随机初始化通常更可靠；若 $K$ 较大，应更加注意初始化质量。
\end{tipbox}
\begin{examplebox}[选择 $K$]
肘部法绘制不同 $K$ 对应的代价，寻找下降速度明显变慢的位置。若没有清晰的肘部，也可以根据业务目标，例如压缩后希望保留多少类、库存需要多少组，选择具有实际意义的 $K$。
\end{examplebox}

\section{主成分分析 PCA}
\subsection{动机：压缩与可视化}
\begin{definitionbox}[低维表示]
若数据位于高维空间但大部分变化集中在少数方向，可以用 $k<n$ 个方向近似表示样本。PCA 选择投影后重建误差较小的方向，而不是简单删除某些坐标。
\end{definitionbox}
\begin{keybox}[PCA 与线性回归的区别]
PCA 是无监督的，不使用预测标签 $y$，目标是最小化正交投影的重建误差；线性回归则选择能预测 $y$ 的方向。不要为了压缩输入而直接运行线性回归或把 PCA 方向当作监督目标。
\end{keybox}
\subsection{PCA 算法}
\begin{definitionbox}[标准步骤]
对每个特征做均值归一化（必要时再除以标准差），得到 $x^{(i)}$。计算协方差矩阵
\[
\Sigma=\frac1mX^TX.
\]
对 $\Sigma$ 做特征值分解或奇异值分解，选取最大的 $k$ 个特征值对应的单位特征向量组成矩阵 $U_{reduce}\in\mathbb R^{n\times k}$，然后投影
\[
z^{(i)}=U_{reduce}^Tx^{(i)}.
\]
\end{definitionbox}
\begin{examplebox}[二维到一维]
若二维样本主要沿一条斜线分布，PCA 会选择这条方向作为第一主成分，把点投影成一个标量 $z$。这个标量保留了数据中最大的方差，适合可视化或降低后续模型的输入维度。
\end{examplebox}
\subsection{选择主成分数量}
\begin{tipbox}[保留方差比例]
选择最小的 $k$，使投影后保留的方差比例满足
\[
\frac{\sum_{j=1}^{k}\lambda_j}{\sum_{j=1}^{n}\lambda_j}\geq0.99
\]
或任务需要的其他阈值。也可以直接比较重建误差。$k$ 不能只凭“看起来方便”决定。
\end{tipbox}
\begin{definitionbox}[重建与信息损失]
低维向量可以近似重建为
\[
\hat x^{(i)}=U_{reduce}z^{(i)}.
\]
重建误差 $\|x^{(i)}-\hat x^{(i)}\|^2$ 越小，说明压缩丢失的信息越少。
\end{definitionbox}
\begin{mistakebox}[PCA 的使用建议]
\begin{itemize}
  \item PCA 应只在训练集上拟合均值、尺度和主成分，再应用到验证集和测试集；
  \item 不要把 PCA 当作解决过拟合的默认方法，先用验证误差确认它确实有帮助；
  \item 对监督学习任务，若原始特征数量并不大，直接用正则化往往更易解释；
  \item 特征尺度差异很大时不做缩放，会让 PCA 主要反映量纲而非真实变化。
\end{itemize}
\end{mistakebox}

\section{本周知识脉络}
\begin{keybox}[聚类与降维]
K-均值通过“分配—移动”迭代最小化簇内距离；PCA 通过最大化保留方差、最小化重建误差寻找低维方向。前者输出组结构，后者输出连续坐标。二者都需要处理尺度、初始化或维度选择，并用下游目标验证实际价值。
\end{keybox}
\end{document}
